%! Function Model Construction and Application — Unit 1, Lesson 1.14 — Guided Notes (Answer Key)
\documentclass[10pt]{article}
\usepackage{precalculus-article}
\usepackage{precalculus-key}   % pulls in -boxes; do NOT also load -boxes
\newcommand{\TallMath}[1]{$\displaystyle #1\rule[-1.4em]{0pt}{3.2em}$}

\begin{document}

\pageheader{Unit 1, Lesson 1.14}{Guided Notes --- Answer Key}
\namedateperiod

\vspace{0.08in}

\begin{objectivebox}
By the end of this lesson, I will be able to\ldots
\begin{enumerate}[itemsep=2pt]
  \item Construct a polynomial or piecewise function model using restrictions, transformations,
        or regression. \hfill\textit{(LO 1.14.A)}
  \item Construct a rational function model from an inverse-proportion context. \hfill\textit{(LO 1.14.B)}
  \item Apply a function model to predict values and compute average rates of change
        with appropriate units. \hfill\textit{(LO 1.14.C)}
\end{enumerate}
\end{objectivebox}

\vspace{0.06in}

\begin{vocabbox}
\termblanklong{function model} \ans{A function constructed to represent a real-world or mathematical scenario; used to make predictions.}\\[2pt]
\termblanklong{domain restriction} \ans{A constraint on the input values of a model based on the context (e.g., length must be positive).}\\[2pt]
\termblanklong{inverse proportion} \ans{A relationship where one quantity equals a constant divided by a power of another: $y = k/x^n$.}\\[2pt]
\termblanklong{average rate of change} \ans{$\dfrac{f(b)-f(a)}{b-a}$; the slope of the secant line between $(a,f(a))$ and $(b,f(b))$.}\\[2pt]
\termblanklong{regression} \ans{A technology-assisted process for finding the best-fit curve of a given type through a data set.}
\end{vocabbox}

\vspace{0.06in}

\begin{hookbox}
A farmer has exactly 80 feet of fencing and wants to build the biggest rectangular pen possible.
One side of the pen sits against a barn, so no fencing is needed on that side.
\textbf{Can a function tell us how to maximize area?}

\medskip
Let $w = $ width (in feet) of the pen. Draw a quick label: two sides of width $w$ and one side
of length $\ell$, with the barn on the fourth side.

\medskip
\textbf{Question:} What equation relates $w$ and $\ell$ to the 80 feet of fencing?

\vspace{0.06in}
\ansline{$2w + \ell = 80$}
\end{hookbox}

\vspace{0.06in}

\begin{notesbox}{Section 1 --- Building a Model from Restrictions (LO 1.14.A)}

\textbf{Setup.} A rectangular pen has one side against a barn (no fencing needed there).
Total fencing available: 80 ft. Two widths and one length must use all 80 ft.

\medskip
\textbf{Step 1.} Write the fencing constraint: \quad $2w + \ell = $ \ans{80}

\medskip
\textbf{Step 2.} Solve for $\ell$ in terms of $w$: \quad $\ell = $ \ans{$80 - 2w$}

\medskip
\textbf{Step 3.} Area of the pen: $A = w \cdot \ell$.  Substitute your expression for $\ell$:

\[
  A(w) = w \cdot \bigl({\color{keyred}\mathbf{80 - 2w}}\bigr) = {\color{keyred}\mathbf{-2w^2 + 80w}}
\]

\medskip
\textbf{Step 4.} Identify the \textbf{domain restriction}.
Width must be positive and fencing constraint must hold:

Domain: $w \in ($ \ans{0} $,$ \ans{40} $)$

\medskip
\textbf{Step 5.} What is the shape of $A(w)$? Circle one: \quad \textbf{Linear} \quad \ans{Quadratic} \quad \textbf{Cubic}

Because the leading coefficient is \ans{$-2$}, the parabola opens \ans{downward}.

\medskip
\textbf{Step 6.} The maximum area occurs at the vertex. For $A(w) = aw^2 + bw$:

$w_{\max} = -\dfrac{b}{2a} = {}$ \ans{$-\dfrac{80}{2(-2)} = 20$ ft} $\quad\Longrightarrow\quad A(w_{\max}) = {}$ \ans{$-2(400)+80(20) = 800$} ft$^2$

\begin{teachernote}
$A(w) = -2w^2 + 80w$; vertex at $w = 20$ ft, $A(20) = 800$ ft$^2$.
Domain: $0 < w < 40$.
\end{teachernote}

\end{notesbox}

\vspace{0.06in}

\begin{notesbox}{Section 2 --- Rational Function Model: Inverse Proportion (LO 1.14.B)}

\textbf{Context.} Two electrically charged objects exert a force on each other.
Experiments show the force is \textbf{inversely proportional to the square of the distance}
between them.

\medskip
This means: \quad $F \propto \dfrac{1}{d^2}$, so we write $F = \dfrac{k}{d^2}$ for some constant $k > 0$.

\medskip
\textbf{Model:} $\quad r(d) = \dfrac{k}{d^2}$

\medskip
\textbf{Domain restriction:} $d$ = distance in meters, so $d$ \ans{$>$} $0$.

\medskip
\textbf{End behavior} (as inputs grow large or approach zero):

\medskip
\renewcommand{\arraystretch}{1.6}
\begin{tabular}{|c|c|}
\hline
As $d \to 0^+$ & $r(d) \to $ \ans{$+\infty$} (force becomes very \ans{large}) \\
\hline
As $d \to \infty$ & $r(d) \to $ \ans{$0$} (force becomes very \ans{small/negligible}) \\
\hline
\end{tabular}

\medskip
\textbf{Example.} Suppose $k = 900$ (units: N$\cdot$m$^2$). Complete the table:

\medskip
\renewcommand{\arraystretch}{1.4}
\begin{tabular}{|c|c|c|c|c|}
\hline
$d$ (m) & 1 & 3 & 5 & 10 \\
\hline
$r(d) = 900/d^2$ & \ans{900} & \ans{100} & \ans{36} & \ans{9} \\
\hline
\end{tabular}

\end{notesbox}

\vspace{0.06in}

\begin{notesbox}{Section 3 --- Applying the Model (LO 1.14.C)}

\textbf{Return to the pen model:} $A(w) = -2w^2 + 80w$ on $(0, 40)$.

\medskip
\textbf{(a) Evaluate at specific inputs.}

$A(10) = -2({\color{keyred}\mathbf{10}})^2 + 80({\color{keyred}\mathbf{10}}) = -200 + 800 = {}$ \ans{600} ft$^2$

\medskip
$A(20) = -2({\color{keyred}\mathbf{20}})^2 + 80({\color{keyred}\mathbf{20}}) = -800 + 1600 = {}$ \ans{800} ft$^2$

\medskip
\textbf{(b) Average rate of change of area on $[10, 20]$.}

$\text{ARC} = \dfrac{A(20) - A(10)}{20 - 10} = \dfrac{{\color{keyred}\mathbf{200}}}{{\color{keyred}\mathbf{10}}} = {}$ \ans{20}

Units: \ans{ft$^2$ per ft (square feet per foot of width)}

\medskip
\textbf{Interpret:} On average, for each additional foot of width between $w = 10$ ft and $w = 20$ ft,
the area \ans{increases} by \ans{20} square feet.

\medskip
\textbf{(c) Predict the width at maximum area.}

From Section 1, the maximum occurs at $w = $ \ans{20} ft with area $= $ \ans{800} ft$^2$.

The length at that width is $\ell = 80 - 2({\color{keyred}\mathbf{20}}) = {}$ \ans{40} ft.

\medskip
\textbf{Conclusion:} The optimal pen is \ans{20} ft wide and \ans{40} ft long.

\end{notesbox}

\end{document}
